Every policy receives the same observable state; capacity regimes scale
total capacity to 70\%/90\%/110\% of mean daily test demand. Decision
loss is cumulative priority-weighted unserved workload; we report each
policy's position between the forecast-free uniform floor and the
perfect-information myopic reference (``share of gap closed''). Four
findings emerge.

\paragraph{Under scarcity, the allocation rule matters more than the
forecaster.}
Proportional allocation driven by the \emph{best} forecaster closes only
56.4\% (New York), 67.2\% (Chicago), 91.4\% (San Francisco), and 64.9\%
(Austin) of the gap in the scarce regime --- while the expected-value
greedy policy driven by the \emph{naive trailing-mean} forecast closes
80.0--95.6\%. Upgrading the policy beats upgrading the forecast. The
two levers are complementary: greedy with the best forecaster reaches
93.7--98.7\% under scarcity.

\paragraph{The decision value of forecast quality is regime-dependent
(H2).}
Holding the greedy policy fixed, replacing the naive forecast with the
best forecast closes an additional 7.8 (New York), 13.7 (Chicago), 2.6
(San Francisco), and 9.4 (Austin) percentage points of the gap under
scarcity, but only 0.6, 8.5, 0.8, and 10.2 points under generous
capacity. In three of four cities the marginal value of accuracy shrinks
substantially as capacity grows; Austin, whose demand is smallest and
spikiest, is the exception and is reported as such.

\paragraph{Accuracy rankings and decision rankings broadly agree --- and
still disagree exactly where it matters (H3).}
Across the eleven point-forecast configurations, Spearman correlation
between MAE ranking and decision-loss ranking is high in every
city-regime pair ($\rho = 0.73$--$0.98$) but never 1. The divergence
concentrates at the top: selecting the deployment model by validation
decision loss instead of validation MAE changes the selected model in
\textbf{7 of 12} city-regime pairs --- always flipping between the local
and pooled variants of the same LightGBM configuration, whose MAE
difference is small but whose allocation behaviour differs. A cheap,
optimization-free form of decision-focused learning --- selecting among
accuracy-trained models by downstream loss --- is therefore advisable
whenever the deployment regime is known.

\paragraph{Forecast uncertainty has operational value under scarcity
(H4).}
Replacing the point forecast with the quantile forecast inside the same
greedy policy improves the scarce-regime gap closure in all four cities
(New York 95.9 $\rightarrow$ 97.0, Chicago 93.7 $\rightarrow$ 96.9, San
Francisco 98.2 $\rightarrow$ 99.8, Austin 98.7 $\rightarrow$ 99.3) ---
despite the under-dispersion documented in Section~6. Under moderate and
generous capacity the comparison is mixed ($\pm$0.3--1.0 points), as
hypothesized: uncertainty matters most when capacity cannot cover the
upper tail everywhere. The equal-weights sensitivity sharpens this
picture: with no priority gradient to exploit, the point-forecast greedy
policy falls behind even proportional allocation in New York, while the
quantile policy stays within 0.6--4.2\% of the perfect-information
reference in every city --- distributional information substitutes for
the structure that priorities otherwise provide.

\begin{figure}[t]\centering
\includegraphics[width=0.95\linewidth]{fig3_accuracy_vs_decision.pdf}
\caption{Forecast accuracy (test MAE, $x$) versus decision value (share
of uniform-to-reference gap closed, $y$) for the eleven point-forecast
configurations under the expected-value greedy policy. Lower MAE
generally buys more decision value, but the slope and spread depend on
city and capacity regime.}
\label{fig:acc-dec}
\end{figure}
